%!TEX root = ../main.tex

\mainmatter
\chapter{Introduction}
\label{introduction}


The numerical evaluation of high-dimensional integrals is a recurring challenge in applied mathematics, computational physics and modern data science. In many cases, such integrals cannot be computed in closed form and must be approximated numerically. The Monte Carlo (MC) method has long been the standard approach, valued for its simplicity and dimension-independent convergence rate. However, its stochastic nature leads to slow convergence, with an error rate of order $O(N^{-1/2})$, which quickly becomes a bottleneck in applications that demand high accuracy or involve costly function evaluations.  

Quasi-Monte Carlo (QMC) methods offer a deterministic alternative. By replacing random samples with carefully constructed low-discrepancy sequences, QMC aims to distribute sampling points more uniformly across the integration domain. Under suitable smoothness assumptions, this improved equidistribution yields substantially faster convergence rates. The theoretical framework, culminating in the Koksma--Hlawka inequality, provides deterministic error bounds and connects discrepancy of the sampling set with the variation of the integrand. QMC methods have therefore attracted increasing attention both in numerical analysis and in practical fields where high-dimensional integration is unavoidable.

The central objective of this thesis is to investigate the theoretical properties and practical benefits of QMC methods in comparison to classical Monte Carlo sampling. The thesis explores how QMC techniques can provide advantages in high-dimensional problems and illustrates their effectiveness in two distinct application areas: the training of deep neural networks and the simulation of photon transport in X-ray imaging. By combining rigorous mathematical foundations with applied case studies, the work aims to highlight both the potential and the limitations of QMC approaches in contemporary computational mathematics.

This objective is pursued in three parts:

\begin{itemize}
    \item \textbf{Part I} develops the mathematical foundations of QMC methods for integral approximation. Central notions such as uniform distribution modulo one, discrepancy and Hardy--Krause variation are introduced. The Koksma--Hlawka inequality is then derived, offering a deterministic error bound that motivates the later applications. 
    
    \item \textbf{Part II} applies QMC sampling to the training of deep neural networks (DNNs). Recent work by \citeauthor{mishra2021enhancing} \cite{mishra2021enhancing} suggests that the generalization error of DNNs can be analyzed in analogy to integration error bounds, with low-discrepancy sequences reducing the generalization gap compared to random sampling. A central theorem in this context is revisited and proven under stricter assumptions, namely that the target function must exhibit $L^1$-continuity of all first-order mixed partial derivatives. Beyond the theoretical contribution, the approach is evaluated empirically on two representative problems: a high-dimensional sum-of-sines function and a surrogate for projectile motion. The results provide evidence that QMC sampling is a promising choice for generating training data in many-query problems.
    
    \item \textbf{Part III} explores a second, independent application: X-ray photon imaging simulation. Here, QMC methods are employed in the context of Monte Carlo-based photon transport, a key tool for modeling scatter in computed tomography (CT). This part begins with an introduction to the physics of X-ray generation, attenuation and scattering. Inspired by the promising results of the gQMCFFD algorithm in \cite{qmcXray2023}, a simulation framework is then developed, incorporating QMC-based sampling strategies into the photon generation and propagation steps. The framework is evaluated on different phantom geometries, showing measurable improvements in accuracy and noise reduction compared to standard Monte Carlo simulations. 
\end{itemize}

It is important to emphasize that Parts II and III are independent of one another: both address applications of QMC methods, but in distinct problem classes. This dual perspective is intended not as a limitation but as a strength, showcasing the breadth of QMC techniques and their potential across disciplines.

\section*{Structure of the Thesis}

The remainder of the thesis is structured as follows:

\begin{itemize}
    \item \textbf{Part~\ref{part1} (Chapters \ref{chapter1}--\ref{chapter3})} introduces Monte Carlo and Quasi-Monte Carlo integration. Chapter~\ref{chapter1} motivates the problem of high-dimensional integration and explains the principle and convergence of MC methods. Chapter~\ref{chapter2} presents discrepancy theory and low-discrepancy sequences, with a focus on the Halton and Sobol' sequences. Chapter~\ref{chapter3} introduces notions of function variation and derives the Koksma--Hlawka inequality.
    
    \item \textbf{Part~\ref{part2} (Chapters \ref{chapter4}--\ref{chapter7})} applies QMC methods to neural network training. Chapter~\ref{chapter4} motivates many-query problems in scientific computing and formulates the problem setting. Chapter~\ref{chapter5} introduces the theoretical foundations of deep neural networks, including function approximation and optimization. Chapter~\ref{chapter6} develops a QMC-based training algorithm and provides theoretical error estimates. Chapter~\ref{chapter7} presents empirical evaluations and discusses the results.
    
    \item \textbf{Part~\ref{part3} (Chapters \ref{chapter8}--\ref{chapter11})} turns to photon transport simulation. Chapter~\ref{chapter8} motivates the problem of scatter in X-ray imaging. Chapter~\ref{chapter9} introduces the underlying physical models of X-ray generation, attenuation and scattering. Chapter~\ref{chapter10} presents the simulation algorithm, combining these physical laws with QMC-based sampling strategies. Chapter~\ref{chapter11} evaluates the framework on different phantoms and discusses results, limitations and future work.
\end{itemize}
